\section{Decision \& Optimization Engine}
\label{sec:decision_optimization}

\subsection{Bridging Predictions to Grid Dispatch}
\label{subsec:bridging_predictions}

The Decision and Optimization Engine transforms generation forecasts into operational grid recommendations. It identifies upcoming imbalances and evaluates candidate dispatch options while respecting physical and regulatory constraints (\cref{fig:decision_pipeline}).

\begin{figure}[htbp]
    \centering
    \begin{tikzpicture}[
        node distance=0.32cm,
        pipeBox/.style={rectangle, rounded corners=3pt, draw=accentTeal!80, fill=boxBg, text width=7.8cm, align=left, font=\scriptsize, inner sep=3pt},
        arrow/.style={-{Stealth[scale=0.8]}, thick, color=darkSlate}
    ]
        \node[pipeBox] (d1) {\textbf{1. Generation Forecast \& Risk Envelope} \hfill \color{darkSlate!70}Hourly Point + Prediction Bands};
        \node[pipeBox, below=of d1] (d2) {\textbf{2. Surplus / Shortfall Event Detection} \hfill \color{darkSlate!70}Imbalance Windows \& Ramp Rate Alerts};
        \node[pipeBox, below=of d2] (d3) {\textbf{3. Candidate Action Formulation} \hfill \color{darkSlate!70}Storage Dispatch, Import/Export, Backup, Curtailment};
        \node[pipeBox, below=of d3] (d4) {\textbf{4. What-If Scenario Stress Testing} \hfill \color{darkSlate!70}Operator Parameter Perturbations (Weather, SOC)};
        \node[pipeBox, below=of d4] (d5) {\textbf{5. Constraint \& Multi-Criteria Evaluation} \hfill \color{darkSlate!70}Operational Feasibility, Cost, Carbon, Wastage};
        \node[pipeBox, below=of d5] (d6) {\textbf{6. Explainable Recommendation Generation} \hfill \color{darkSlate!70}Structured Advice: Action + Reason + Impact};
        \node[pipeBox, below=of d6, draw=accentAmber!90, fill=accentAmber!10] (d7) {\textbf{7. Human Operator Decision} \hfill \color{darkSlate!70}\textbf{Supervisory Review \& Dispatch Approval}};

        \draw[arrow] (d1) -- (d2);
        \draw[arrow] (d2) -- (d3);
        \draw[arrow] (d3) -- (d4);
        \draw[arrow] (d4) -- (d5);
        \draw[arrow] (d5) -- (d6);
        \draw[arrow] (d6) -- (d7);
    \end{tikzpicture}
    \caption{Decision \& Optimization Engine operational pipeline ending in human operator approval.}
    \label{fig:decision_pipeline}
\end{figure}

\subsection{Operational Pathways \& Multi-Criteria Evaluation}
\label{subsec:operational_pathways}

The engine formulates tailored response sequences based on detected conditions:
\begin{itemize}[leftmargin=1.5em, itemsep=0.2em]
    \item \textbf{Under-Generation (Shortfall):} Prioritizes BESS discharge $\rightarrow$ grid power imports within intertie limits $\rightarrow$ thermal/hydro backup activation.
    \item \textbf{Over-Generation (Surplus):} Prioritizes BESS charging $\rightarrow$ inter-regional power export $\rightarrow$ managed curtailment.
\end{itemize}

Candidate actions are ranked across five criteria: \emph{technical feasibility}, \emph{renewable energy utilization}, \emph{operating cost impact}, \emph{carbon intensity}, and \emph{reliability reserve margins}.

\subsection{Explainable Decision Support}
\label{subsec:explainable_output}

Every recommendation delivered to the control room follows a transparent structure:
\begin{center}
    \fbox{\parbox{0.92\textwidth}{\scriptsize
        \textbf{\color{primaryNavy}RECOMMENDED ACTION:} Dispatch BESS at 18~MW from 17:30 to 19:30.\\
        \textbf{\color{primaryNavy}OPERATIONAL REASON:} Solar generation falls 18~MW below demand; battery SOC (78\%) is sufficient to bridge the deficit.\\
        \textbf{\color{primaryNavy}EXPECTED IMPACT:} Prevents fossil-fuel peaker startup, lowering operating costs and emissions while maintaining frequency stability.
    }}
\end{center}
GridSense AI functions strictly as a decision-support platform; it does not autonomously command physical electrical grid infrastructure.
